\documentclass[10pt,a4paper]{amsart}
\usepackage{amsmath,amssymb,amsthm,graphicx}
% Retained preamble for the Stacks Project chapter snapshot in this directory.
% The chapter's own preamble.tex is a copy of the shared Stacks Project preamble; the block
% below is sliced MECHANICALLY, line for line, from it. Nothing was added, removed or
% redefined; the wrapper class, the xr-hyper/hyperref wiring and the externaldocument list
% are omitted because the assembler supplies the document class.
%
% lines taken from preamble.tex: 3, 5, 9, 11, 12, 156, 157, 158, ... (51 lines total)
%
% package coverage restored from the chapter's own preamble (2026-10-13): xy, xy-2cell

\usepackage[all]{xy}
\xyoption{2cell}
\UseAllTwocells

\newenvironment{reference}{\comment}{\endcomment}
\newenvironment{slogan}{\comment}{\endcomment}
\newenvironment{history}{\comment}{\endcomment}
\theoremstyle{plain}
\newtheorem{theorem}[subsection]{Theorem}
\newtheorem{proposition}[subsection]{Proposition}
\newtheorem{lemma}[subsection]{Lemma}
\theoremstyle{definition}
\newtheorem{definition}[subsection]{Definition}
\newtheorem{example}[subsection]{Example}
\newtheorem{exercise}[subsection]{Exercise}
\newtheorem{situation}[subsection]{Situation}
\theoremstyle{remark}
\newtheorem{remark}[subsection]{Remark}
\newtheorem{remarks}[subsection]{Remarks}
\numberwithin{equation}{subsection}
\def\lim{\mathop{\mathrm{lim}}\nolimits}
\def\colim{\mathop{\mathrm{colim}}\nolimits}
\def\Spec{\mathop{\mathrm{Spec}}}
\def\Hom{\mathop{\mathrm{Hom}}\nolimits}
\def\Ext{\mathop{\mathrm{Ext}}\nolimits}
\def\SheafHom{\mathop{\mathcal{H}\!\mathit{om}}\nolimits}
\def\SheafExt{\mathop{\mathcal{E}\!\mathit{xt}}\nolimits}
\def\Sch{\mathit{Sch}}
\def\Mor{\mathop{\mathrm{Mor}}\nolimits}
\def\Ob{\mathop{\mathrm{Ob}}\nolimits}
\def\Sh{\mathop{\mathit{Sh}}\nolimits}
\def\NL{\mathop{N\!L}\nolimits}
\def\CH{\mathop{\mathrm{CH}}\nolimits}
\def\proetale{{pro\text{-}\acute{e}tale}}
\def\etale{{\acute{e}tale}}
\def\QCoh{\mathit{QCoh}}
\def\Ker{\mathop{\mathrm{Ker}}}
\def\Im{\mathop{\mathrm{Im}}}
\def\Coker{\mathop{\mathrm{Coker}}}
\def\Coim{\mathop{\mathrm{Coim}}}
\DeclareMathSymbol{\boxtimes}{\mathbin}{AMSa}{"02}
\def\QCohstack{\mathcal{QC}\!\mathit{oh}}
\def\Cohstack{\mathcal{C}\!\mathit{oh}}
\def\Spacesstack{\mathcal{S}\!\mathit{paces}}
\def\Quotfunctor{\mathrm{Quot}}
\def\Hilbfunctor{\mathrm{Hilb}}
\def\Curvesstack{\mathcal{C}\!\mathit{urves}}
\def\Polarizedstack{\mathcal{P}\!\mathit{olarized}}
\def\Complexesstack{\mathcal{C}\!\mathit{omplexes}}
\def\Pic{\mathop{\mathrm{Pic}}\nolimits}
\def\Picardstack{\mathcal{P}\!\mathit{ic}}
\def\Picardfunctor{\mathrm{Pic}}
\def\Deformationcategory{\mathcal{D}\!\mathit{ef}}
\usepackage{cleveref}
\crefname{originallemma}{original lemma}{original lemmas}
\Crefname{originallemma}{Original lemma}{Original lemmas}
\makeatletter
\newlabel{descent-lemma-galois-descent}{{0CDR}{source}{}{sourcepointer-1}{}}
\newlabel{descent-lemma-galois-descent@cref}{{[originallemma][0][]0CDR}{[1][0][]source}{}{}{}}
\makeatother
\crefname{originallemma}{original lemma}{original lemmas}
\Crefname{originallemma}{Original lemma}{Original lemmas}
\makeatletter
\newlabel{fields-lemma-infinite-galois-limit}{{0BU2}{source}{}{sourcepointer-2}{}}
\newlabel{fields-lemma-infinite-galois-limit@cref}{{[originallemma][0][]0BU2}{[1][0][]source}{}{}{}}
\makeatother
\crefname{originallemma}{original lemma}{original lemmas}
\Crefname{originallemma}{Original lemma}{Original lemmas}
\makeatletter
\newlabel{limits-lemma-descend-invertible-modules}{{0B8W}{source}{}{sourcepointer-3}{}}
\newlabel{limits-lemma-descend-invertible-modules@cref}{{[originallemma][0][]0B8W}{[1][0][]source}{}{}{}}
\makeatother
\crefname{originallemma}{original lemma}{original lemmas}
\Crefname{originallemma}{Original lemma}{Original lemmas}
\makeatletter
\newlabel{limits-section-limits}{{01YV}{source}{}{sourcepointer-4}{}}
\newlabel{limits-section-limits@cref}{{[originallemma][0][]01YV}{[1][0][]source}{}{}{}}
\makeatother
\crefname{originallemma}{original lemma}{original lemmas}
\Crefname{originallemma}{Original lemma}{Original lemmas}
\makeatletter
\newlabel{varieties-lemma-change-fields-pic}{{0CC5}{source}{}{sourcepointer-5}{}}
\newlabel{varieties-lemma-change-fields-pic@cref}{{[originallemma][0][]0CC5}{[1][0][]source}{}{}{}}
\makeatother
\title{Pic of X is the Galois-invariant Pic of X\_\{k'\}}
\author{The Stacks Project Authors}
\date{stacks-project:pic.tex:local-snapshot-20261012; GNU FDL 1.2 or later}
\begin{document}
\maketitle
\noindent\textit{Editorial note.} Statement span is lines 1558-1568 of sources/stacks-pic/pinned-original-source.tex; proof span is lines 1570-1623, the complete author proof (52 body lines).\par\medskip
\noindent\textit{Editorial note.} Navigation stubs: 5 cited-but-not-included labels are supplied as non-load-bearing pointers with their pinned tags where the snapshot tag index has them.\par\medskip
\noindent\textit{Editorial note.} extra\_packages is empty on purpose: the retained preamble already loads xy with options and a second declaration clashes in the assembler wrapper.\par\medskip
\noindent\textit{Editorial note.} Text is retained verbatim, including the author's own typos and placeholders.\par\medskip
\noindent\textit{Editorial note.} Queue index 5 of queue\_pic.json (pb=2833, stmt 1558-1568, proof 1570-1623).\par\medskip
\noindent\textit{Editorial note.} stacks-pic wave-1 (pi01): fresh adjudication accept (2026-10-14); pool per standing queue hypercovering -> pic (SMC reassigned to Zhen per Chief 08:5x); pic.tex 27 labels/23 with proof judged 21 accept / 2 reject / 4 pb=0 env-dup rows excluded; spans cite-free and reference/slogan/history-free so no patchenv needed.\par\medskip
\noindent\textit{Allowed standard prerequisites (editorial).}\begin{itemize}
\item Explicit prerequisite boundary. The cited inputs are prerequisites; the argument in the retained span is not replaced by any cited result.
\end{itemize}\medskip
\ifdefined\hypertarget\hypertarget{sourcepointer-1}{}\fi
\noindent\textit{Source reference.} Original lemma 0CDR: sources/stacks-pic/pinned-original-source.tex, source lines 1623--1623. Cited by the retained statement or proof but not part of the collected spans; supplied only as a navigation pointer. The argument does not depend on its content.\par\medskip
\ifdefined\hypertarget\hypertarget{sourcepointer-2}{}\fi
\noindent\textit{Source reference.} Original lemma 0BU2: sources/stacks-pic/pinned-original-source.tex, source lines 1623--1623. Cited by the retained statement or proof but not part of the collected spans; supplied only as a navigation pointer. The argument does not depend on its content.\par\medskip
\ifdefined\hypertarget\hypertarget{sourcepointer-3}{}\fi
\noindent\textit{Source reference.} Original lemma 0B8W: sources/stacks-pic/pinned-original-source.tex, source lines 1623--1623. Cited by the retained statement or proof but not part of the collected spans; supplied only as a navigation pointer. The argument does not depend on its content.\par\medskip
\ifdefined\hypertarget\hypertarget{sourcepointer-4}{}\fi
\noindent\textit{Source reference.} Original lemma 01YV: sources/stacks-pic/pinned-original-source.tex, source lines 1623--1623. Cited by the retained statement or proof but not part of the collected spans; supplied only as a navigation pointer. The argument does not depend on its content.\par\medskip
\ifdefined\hypertarget\hypertarget{sourcepointer-5}{}\fi
\noindent\textit{Source reference.} Original lemma 0CC5: sources/stacks-pic/pinned-original-source.tex, source lines 1623--1623. Cited by the retained statement or proof but not part of the collected spans; supplied only as a navigation pointer. The argument does not depend on its content.\par\medskip
\begin{lemma}
\label{lemma-pic-descends}
Let $k$ be a field. Let $X$ be a quasi-compact and quasi-separated scheme
over $k$ with $H^0(X, \mathcal{O}_X) = k$. If $X$ has a $k$-rational point,
then for any Galois extension $k'/k$ we have
$$
\Pic(X) = \Pic(X_{k'})^{\text{Gal}(k'/k)}
$$
Moreover the action of $\text{Gal}(k'/k)$ on $\Pic(X_{k'})$
is continuous.
\end{lemma}

\begin{proof}
Since $\text{Gal}(k'/k) = \text{Aut}(k'/k)$ it acts (from the right)
on $\Spec(k')$, hence it acts (from the right) on
$X_{k'} = X \times_{\Spec(k)} \Spec(k')$, and since $\Pic(-)$
is a contravariant functor, it acts (from the left) on $\Pic(X_{k'})$.
If $k'/k$ is an infinite Galois extension, then we write
$k' = \colim k'_\lambda$ as a filtered colimit of finite Galois
extensions, see Fields, Lemma \ref{fields-lemma-infinite-galois-limit}.
Then $X_{k'} = \lim X_{k_\lambda}$
(as in Limits, Section \ref{limits-section-limits})
and we obtain
$$
\Pic(X_{k'}) = \colim \Pic(X_{k_\lambda})
$$
by Limits, Lemma \ref{limits-lemma-descend-invertible-modules}.
Moreover, the transition maps in this system of abelian groups
are injective by Varieties, Lemma \ref{varieties-lemma-change-fields-pic}.
It follows that every element of $\Pic(X_{k'})$ is fixed by
one of the open subgroups $\text{Gal}(k'/k_\lambda)$, which exactly
means that the action is continuous. Injectivity of the transition maps
implies that it suffices to prove the statement on fixed points in the
case that $k'/k$ is finite Galois.

\medskip\noindent
Assume $k'/k$ is finite Galois with Galois group $G = \text{Gal}(k'/k)$.
Let $\mathcal{L}$ be an element of $\Pic(X_{k'})$ fixed by $G$.
We will use Galois descent (Descent, Lemma \ref{descent-lemma-galois-descent})
to prove that $\mathcal{L}$ is the pullback of an invertible sheaf on $X$.
Recall that $f_\sigma = \text{id}_X \times \Spec(\sigma) : X_{k'} \to X_{k'}$
and that $\sigma$ acts on $\Pic(X_{k'})$ by pulling back by $f_\sigma$.
Hence for each $\sigma \in G$ we can choose an isomorphism
$\varphi_\sigma : \mathcal{L} \to f_\sigma^*\mathcal{L}$
because $\mathcal{L}$ is a fixed by the $G$-action.
The trouble is that we don't know if we can choose
$\varphi_\sigma$ such that the cocycle condition
$\varphi_{\sigma\tau} = f_\sigma^*\varphi_\tau \circ \varphi_\sigma$
holds. To see that this is possible we use that $X$ has a $k$-rational point
$x \in X(k)$. Of course, $x$ similarly determines a $k'$-rational point
$x' \in X_{k'}$ which is fixed by $f_\sigma$ for all $\sigma$.
Pick a nonzero element $s$ in the fibre of $\mathcal{L}$ at $x'$;
the fibre is the $1$-dimensional $k' = \kappa(x')$-vector space
$$
\mathcal{L}_{x'} \otimes_{\mathcal{O}_{X_{k'}, x'}} \kappa(x').
$$
Then $f_\sigma^*s$ is a nonzero element of the fibre of
$f_\sigma^*\mathcal{L}$ at $x'$. Since we can multiply $\varphi_\sigma$
by an element of $(k')^*$ we may assume that $\varphi_\sigma$ sends
$s$ to $f_\sigma^*s$. Then we see that both $\varphi_{\sigma\tau}$ and
$f_\sigma^*\varphi_\tau \circ \varphi_\sigma$
send $s$ to $f_{\sigma\tau}^*s = f_\tau^*f_\sigma^*s$.
Since $H^0(X_{k'}, \mathcal{O}_{X_{k'}}) = k'$ these two isomorphisms
have to be the same (as one is a global unit times the other
and they agree in $x'$) and the proof is complete.
\end{proof}

\end{document}
